% GENERATED FILE. Edit canonical lessons, then run npm run generate:books.
% canonical-source-hash: fnv1a64:f9aa50b1496bb45c
% canonical-lessons: ZH-C239-jiankang, ZH-C239-jieyan, ZH-C239-hejiu, ZH-C239-huxi, ZH-C239-tijian

\chapter{Healthy Habits}
\label{ch:zh-healthy-habits}
\hlchaptermodality{239}

\begin{chapteropening}
\textbf{By the end of this chapter:} \emph{I can talk about staying healthy, giving up smoking and drinking, breathing and having a check-up.}

Complete the last lesson of chapter 239: I can talk about staying healthy, giving up smoking and drinking, breathing and having a check-up.
\end{chapteropening}

\section[jiànkāng]{jiànkāng — healthy; health}

\label{lesson:ZH-C239-jiankang}



\begin{quote}
Before the new one: say the Mandarin for the open sea, then the Mandarin for a rainbow.
\end{quote}

\subsection*{You'll want to know: jiànkāng}
\emph{jiànkāng} — \textquotedblleft{}healthy; health\textquotedblright{}.

\begin{quote}
\emph{tā měitiān pǎobù hěn jiànkāng} — He runs every day and he's very fit.
\end{quote}

Its characters come later; for now it is a word for the ear.

\begin{cousinweb}
\emph{jiànkāng} is built from \emph{jiàn} 'strong' and \emph{kāng} 'at ease', both inherited from Middle Chinese, \emph{gjonH} and \emph{khang} (Guangyun, in Baxter's transcription), which glosses the second 'in harmony' and 'happy'.
\end{cousinweb}

\begin{grammarlens}[title={What you've built}]
One more word for saying what I do.
\end{grammarlens}

\subsection*{Your turn}
\begin{itemize}
\raggedright
  \item \emph{Say it:} \emph{jiànkāng}
  \item \emph{Say it:} \emph{jiànkāng}, once more
  \item \emph{Say it:} \emph{tā měitiān pǎobù hěn jiànkāng} — He runs every day and he's very fit.
  \item \emph{Recall:} say the Mandarin for the open sea, then the Mandarin for a rainbow, then say \emph{jiànkāng} again
\end{itemize}

\begin{tcolorbox}[breakable,colback=teal!4,colframe=teal!35!black,title={Before you move on}]
What does jiànkāng mean? (\textquotedblleft{}Healthy; health\textquotedblright{}.) Say it once more.
\end{tcolorbox}
\section[jièyān]{jièyān — to give up smoking}

\label{lesson:ZH-C239-jieyan}



\begin{quote}
Before the new one: say the Mandarin for a rainbow, then the Mandarin for healthy.
\end{quote}

\subsection*{You'll want to know: jièyān}
\emph{jièyān} — \textquotedblleft{}to give up smoking\textquotedblright{}.

\begin{quote}
\emph{wèile shēntǐ wǒ bàba jièyān le} — My dad gave up smoking for his health.
\end{quote}

Its characters come later; for now it is a word for the ear.

\begin{cousinweb}
\emph{jièyān} is built from \emph{jiè} 'to guard against' and \emph{yān} 'smoke'. The second is inherited from Middle Chinese \emph{'en} (Guangyun, in Baxter's transcription), glossed 'the breath of fire'; the dictionary glosses the first, \emph{keajH} there, 'to be careful' and 'to warn'.
\end{cousinweb}

\begin{grammarlens}[title={What you've built}]
One more word for saying what I do.
\end{grammarlens}

\subsection*{Your turn}
\begin{itemize}
\raggedright
  \item \emph{Say it:} \emph{jièyān}
  \item \emph{Say it:} \emph{jièyān}, once more
  \item \emph{Say it:} \emph{wèile shēntǐ wǒ bàba jièyān le} — My dad gave up smoking for his health.
  \item \emph{Recall:} say the Mandarin for a rainbow, then the Mandarin for healthy, then say \emph{jièyān} again
\end{itemize}

\begin{tcolorbox}[breakable,colback=teal!4,colframe=teal!35!black,title={Before you move on}]
What does jièyān mean? (\textquotedblleft{}To give up smoking\textquotedblright{}.) Say it once more.
\end{tcolorbox}
\section[hējiǔ]{hējiǔ — to drink alcohol}

\label{lesson:ZH-C239-hejiu}



\begin{quote}
Before the new one: say the Mandarin for healthy, then the Mandarin for to give up smoking.
\end{quote}

\subsection*{You'll want to know: hējiǔ}
\emph{hējiǔ} — \textquotedblleft{}to drink alcohol\textquotedblright{}.

\begin{quote}
\emph{tā bù xīyān yě bù hējiǔ} — He doesn't smoke or drink.
\end{quote}

Its characters come later; for now it is a word for the ear.

\begin{cousinweb}
\emph{hējiǔ} is built from \zh{喝} 'to drink' and \emph{jiǔ} 'alcohol', both inherited from Middle Chinese, \emph{xat} and \emph{tsjuwX} (Guangyun, in Baxter's transcription); the dictionary glosses the first 'to shout', so how it came to mean drinking is not settled.
\end{cousinweb}

\begin{grammarlens}[title={What you've built}]
One more word for saying what I do.
\end{grammarlens}

\subsection*{Your turn}
\begin{itemize}
\raggedright
  \item \emph{Say it:} \emph{hējiǔ}
  \item \emph{Say it:} \emph{hējiǔ}, once more
  \item \emph{Say it:} \emph{tā bù xīyān yě bù hējiǔ} — He doesn't smoke or drink.
  \item \emph{Recall:} say the Mandarin for healthy, then the Mandarin for to give up smoking, then say \emph{hējiǔ} again
\end{itemize}

\begin{tcolorbox}[breakable,colback=teal!4,colframe=teal!35!black,title={Before you move on}]
What does hējiǔ mean? (\textquotedblleft{}To drink alcohol\textquotedblright{}.) Say it once more.
\end{tcolorbox}
\section[hūxī]{hūxī — to breathe}

\label{lesson:ZH-C239-huxi}



\begin{quote}
Before the new one: say the Mandarin for to give up smoking, then the Mandarin for to drink alcohol.
\end{quote}

\subsection*{You'll want to know: hūxī}
\emph{hūxī} — \textquotedblleft{}to breathe\textquotedblright{}.

\begin{quote}
\emph{tā gǎnmào le bù néng yòng bízi hūxī} — He's got a cold and can't breathe through his nose.
\end{quote}

Its characters come later; for now it is a word for the ear.

\begin{cousinweb}
\emph{hūxī} is built from \emph{hū} 'to breathe out' and \emph{xī} 'to breathe in', both inherited from Middle Chinese, \emph{xu} and \emph{xip} (Guangyun, in Baxter's transcription), which glosses the first 'to call' and, after an older book, 'the outward breath', and the second 'the inward breath'.
\end{cousinweb}

\begin{grammarlens}[title={What you've built}]
One more word for saying what I do.
\end{grammarlens}

\subsection*{Your turn}
\begin{itemize}
\raggedright
  \item \emph{Say it:} \emph{hūxī}
  \item \emph{Say it:} \emph{hūxī}, once more
  \item \emph{Say it:} \emph{tā gǎnmào le bù néng yòng bízi hūxī} — He's got a cold and can't breathe through his nose.
  \item \emph{Recall:} say the Mandarin for to give up smoking, then the Mandarin for to drink alcohol, then say \emph{hūxī} again
\end{itemize}

\begin{tcolorbox}[breakable,colback=teal!4,colframe=teal!35!black,title={Before you move on}]
What does hūxī mean? (\textquotedblleft{}To breathe\textquotedblright{}.) Say it once more.
\end{tcolorbox}
\section[tǐjiǎn]{tǐjiǎn — medical check-up}

\label{lesson:ZH-C239-tijian}



\begin{quote}
Before the new one: say the Mandarin for to drink alcohol, then the Mandarin for to breathe.
\end{quote}

\subsection*{You'll want to know: tǐjiǎn}
\emph{tǐjiǎn} — \textquotedblleft{}medical check-up\textquotedblright{}.

\begin{quote}
\emph{wǒ měi nián zuò yí cì tǐjiǎn} — I have a check-up once a year.
\end{quote}

Its characters come later; for now it is a word for the ear.

Tone: tǐ and jiǎn are both third tones, so the first rises and it sounds like tíjiǎn.

\begin{cousinweb}
\emph{tǐjiǎn} is built from \zh{体} 'body' and \emph{jiǎn} 'to check', both inherited from Middle Chinese, \emph{thejX} and \emph{kjemX} (Guangyun, in Baxter's transcription).
\end{cousinweb}

\begin{grammarlens}[title={What you've built}]
One more word for saying what I do.
\end{grammarlens}

\subsection*{Your turn}
\begin{itemize}
\raggedright
  \item \emph{Say it:} \emph{tǐjiǎn}
  \item \emph{Say it:} \emph{tǐjiǎn}, once more
  \item \emph{Say it:} \emph{wǒ měi nián zuò yí cì tǐjiǎn} — I have a check-up once a year.
  \item \emph{Recall:} say the Mandarin for to drink alcohol, then the Mandarin for to breathe, then say \emph{tǐjiǎn} again
\end{itemize}

\begin{tcolorbox}[breakable,colback=teal!4,colframe=teal!35!black,title={Before you move on}]
What does tǐjiǎn mean? (\textquotedblleft{}Medical check-up\textquotedblright{}.) Say it once more.
\end{tcolorbox}
